\begin{conf-abstract}
{Zimbabwe-Headwater Catchment Response to Extreme Rainfall Activities-Meeting on sampling of hydrological extreme events}
{Webster	Gumindoga}
{University of Zimbabwe}
{Extreme rainfall events are becoming increasingly frequent across Zimbabwe, yet their effects on headwater catchment hydrology remain poorly quantified. Understanding rainfall–runoff–groundwater interactions under such events is essential for water resources planning, particularly in convective rainfall-dominated systems. This study assessed the hydrological and isotopic response of three representative headwater catchments—Domboshava (Mazowe), Marondera (Manyame), and the University of Zimbabwe (Gwebi)—during the 2023/2024 wet season. Stable isotopes of oxygen (\ce{\delta^18O}) and hydrogen (\ce{\delta^2H}) were analysed in rainfall, surface water, and groundwater, while Radon-222 (\ce{^222Rn}) was used as a tracer of subsurface flow dynamics. Sampling was conducted at event and intra-event scales, with rainfall collected up to three times daily and groundwater monitored through logging wells. Early-season rainfall at the University of Zimbabwe site exhibited isotopic enrichment (\ce{\delta^18O} = 3.06\permil; \ce{\delta^2H} = 17.48\permil) and low d-excess (-7\permil), indicating subcloud evaporation and limited recharge potential. Intra-event sampling revealed “$\Lambda$-shaped” isotopic profiles, reflecting the temporal variability within convective storm systems. Distinct isotopic signatures were observed among rainfall, surface water, and groundwater, underscoring episodic recharge patterns. Variability in \ce{^222Rn} concentrations indicated contrasting subsurface flow paths and groundwater residence times across sites. The study highlights the heterogeneity of recharge processes in Zimbabwean headwater catchments and the influence of convective rainfall on water partitioning. Data limitations due to manual sampling and equipment gaps remain key challenges. Future work will focus on improved monitoring networks, recharge zone mapping, and long-term isotope-hydrology integration for water resource resilience under extreme rainfall conditions. Keywords: Stable isotopes, Convective rainfall, Groundwater recharge, Headwater catchment, Radon-222, Zimbabwe.
}
\end{conf-abstract}


% Extreme rainfall events are becoming increasingly frequent across Zimbabwe, yet their effects on headwater catchment hydrology remain poorly quantified. Understanding rainfall–runoff–groundwater interactions under such events is essential for water resources planning, particularly in convective rainfall-dominated systems. This study assessed the hydrological and isotopic response of three representative headwater catchments—Domboshava (Mazowe), Marondera (Manyame), and the University of Zimbabwe (Gwebi)—during the 2023/2024 wet season. Stable isotopes of oxygen (\ce{\delta^18O}) and hydrogen (δ²H) were analysed in rainfall, surface water, and groundwater, while Radon-222 (²²²Rn) was used as a tracer of subsurface flow dynamics. Sampling was conducted at event and intra-event scales, with rainfall collected up to three times daily and groundwater monitored through logging wells. Early-season rainfall at the University of Zimbabwe site exhibited isotopic enrichment (δ¹⁸O = 3.06‰; δ²H = 17.48‰) and low d-excess (−7‰), indicating subcloud evaporation and limited recharge potential. Intra-event sampling revealed “ʌ-shaped” isotopic profiles, reflecting the temporal variability within convective storm systems. Distinct isotopic signatures were observed among rainfall, surface water, and groundwater, underscoring episodic recharge patterns. Variability in ²²²Rn concentrations indicated contrasting subsurface flow paths and groundwater residence times across sites. The study highlights the heterogeneity of recharge processes in Zimbabwean headwater catchments and the influence of convective rainfall on water partitioning. Data limitations due to manual sampling and equipment gaps remain key challenges. Future work will focus on improved monitoring networks, recharge zone mapping, and long-term isotope-hydrology integration for water resource resilience under extreme rainfall conditions. Keywords: Stable isotopes, Convective rainfall, Groundwater recharge, Headwater catchment, Radon-222, Zimbabwe.
